\documentclass[11pt]{article}

\usepackage[T1]{fontenc}
\usepackage{amsmath,amssymb,amsthm}
\usepackage[margin=1in]{geometry}
\usepackage{microtype}
\usepackage[hidelinks]{hyperref}
\hypersetup{
  pdftitle={Direct Scalene Triangular Facets of Convex Polyhedra Must Be Acute},
  pdfauthor={Xiaojian Zeng},
  pdfsubject={Convex polytopes with orientation-preserving congruent triangular facets}
}

\newtheorem{theorem}{Theorem}
\newtheorem{lemma}[theorem]{Lemma}
\newcommand{\R}{\mathbb{R}}

\title{Direct Scalene Triangular Facets of Convex Polyhedra Must Be Acute}
\author{Xiaojian Zeng\\
  \small PBC School of Finance, Tsinghua University, 100190\\
  \small \texttt{xiaojianzeng@foxmail.com}}
\date{September 10, 2026}

\begin{document}
\maketitle

\begin{abstract}
Fix either planar orientation of a nondegenerate scalene triangle $T$.
We prove that if every facet of a full-dimensional compact convex
three-polytope is congruent to $T$ by an orientation-preserving congruence,
where facet orientation is induced by the outward normal, then $T$ is acute.
Thus no right or obtuse scalene triangle has an orientation-preserving
monohedral realization.  The key local observation is that the side labels
around every oriented vertex link advance by a fixed nontrivial three-cycle.
Consequently every vertex degree is divisible by three; Euler's formula then
forces a degree-three vertex, whose star determines a nondegenerate
disphenoid.  Its elementary coordinate identities imply strict acuteness.
\end{abstract}

\section{Context and statement}

Problem B22 of Croft, Falconer, and Guy asks which triangles can occur as the
common face of a convex polyhedron, both with and without reflected copies
\cite[Problem B22]{CroftFalconerGuy1991}.  Georgiou recently solved the
reflection-allowed existence and minimum-face-count questions
\cite{Georgiou2026}.  For a scalene triangle, however, his eight-face
construction contains four direct and four mirror-image facets.  With cyclic
side order read using the outward normal, his Remark~3 therefore leaves the
orientation-preserving case untouched; his concluding discussion isolates
right and obtuse scalene triangles as the remaining case
\cite[Remark~3 and Section~4]{Georgiou2026}.

We use the following precise convention.  Fix an orientation of a scalene
triangle $T$ with side-length set $\{a,b,c\}$.  Its positive boundary has a
cyclic side-length word.  A triangular facet of a convex three-polytope is
oriented by its outward normal, and it is a \emph{direct copy} of $T$ when its
positive cyclic side-length word is a cyclic rotation of the word for $T$,
not of its reverse.  This agrees exactly with Georgiou's direct-versus-mirror
distinction for scalene facets \cite[Remark~3]{Georgiou2026}.

\begin{theorem}\label{thm:main}
Let $T$ be a nondegenerate scalene Euclidean triangle, with either of its two
planar orientations fixed.  If a full-dimensional compact convex
three-polytope has every facet a direct copy of $T$, then $T$ is acute.
Equivalently, if
\[
  0<a<b<c<a+b, \qquad c^2\geq a^2+b^2,
\]
then no such polytope exists.
\end{theorem}

The proof supplies the orientation-sensitive local step missing from the
reflection-allowed construction and then reduces the problem to the
four-face obstruction.  The latter is the classical disphenoid calculation
studied, for example, by Leech \cite{Leech1950} and used as Lemma~5 by
Georgiou \cite{Georgiou2026}; we reproduce the needed calculation so that the
argument is self-contained.  A bounded literature search through the date of
this manuscript located no later paper proving Theorem~\ref{thm:main} or an
equivalent oriented-link lemma.  This dated negative search is not a claim of
priority.

\section{The oriented-link constraint}

Fix the usual orientation of $\R^3$.  An ordered facet $(p_0,p_1,p_2)$ is
positive when
\[
 ((p_1-p_0)\times(p_2-p_0))\mathbin{\cdot}n_{\mathrm{out}}>0.
\]
Positive boundary orders on two adjacent facets traverse their common edge
in opposite directions.

Let $\psi$ be the next-label permutation in the fixed positive cyclic word
of $T$.  Since $T$ is scalene, $\psi$ is one of the two nonidentity
three-cycles on $\{a,b,c\}$.  Reversing the chosen orientation of $T$
replaces $\psi$ by $\psi^{-1}$.

\begin{lemma}[oriented-link three-cycle]\label{lem:link}
At every vertex $v$, the degree $\deg(v)$ is divisible by three.
More precisely, if $u_0,\ldots,u_{d-1}$ are its cyclic neighbors, indexed so
that every incident facet is positively ordered as
$(v,u_i,u_{i+1})$, and if
\[
 r_i=|vu_i|, \qquad q_i=|u_i u_{i+1}|,
\]
then, with indices modulo $d$,
\begin{equation}\label{eq:transition}
 q_i=\psi(r_i), \qquad r_{i+1}=\psi^2(r_i).
\end{equation}
\end{lemma}

\begin{proof}
The link of $v$ in the triangular boundary sphere is a cycle.  Begin with
one positive incident facet $(v,u_0,u_1)$.  Across $vu_1$, compatibility of
the oriented facet boundaries forces the next positive facet to be
$(v,u_1,u_2)$.  Continuing gives the asserted cyclic indexing.

The positive side word on $(v,u_i,u_{i+1})$ is
$(r_i,q_i,r_{i+1})$.  Directness says that this word is a cyclic rotation of
the fixed word of $T$.  Its three entries are distinct, so its first entry
determines the other two and yields \eqref{eq:transition}.  Closure of the
link gives
\[
 r_0=r_d=\psi^{2d}(r_0).
\]
A nonidentity three-cycle has no fixed label under its first or second power.
Hence $3\mid 2d$, and therefore $3\mid d$.
\end{proof}

Notice that the same proof treats both handedness classes: it uses only that
$\psi$ has order three, not which of the two three-cycles it is.

\section{Euler reduction and the disphenoid}

\begin{lemma}\label{lem:degree3}
The polytope has a vertex of degree three.
\end{lemma}

\begin{proof}
Let $V,E,F$ be the numbers of vertices, edges, and facets.  The boundary is a
triangulated two-sphere, so
\[
 V-E+F=2,\qquad 3F=2E,\qquad
 \sum_v\deg(v)=2E.
\]
Eliminating $E$ and $F$ gives
\begin{equation}\label{eq:curvature}
 \sum_v(6-\deg(v))=12.
\end{equation}
Every vertex has degree at least three, and Lemma~\ref{lem:link} makes every
degree a multiple of three.  If no degree were three, every degree would be
at least six and the left side of \eqref{eq:curvature} would be nonpositive,
a contradiction.
\end{proof}

\begin{lemma}\label{lem:star}
Let $v$ have degree three and cyclic neighbors $u_0,u_1,u_2$.  The four
points $v,u_0,u_1,u_2$ are affinely independent, and their six joining edges
occur in three equal opposite pairs.  In particular, they form a
nondegenerate tetrahedron all four of whose faces are congruent to $T$.
\end{lemma}

\begin{proof}
For the notation of Lemma~\ref{lem:link}, equation
\eqref{eq:transition} gives
\[
 q_{i+1}=\psi(r_{i+1})=\psi^3(r_i)=r_i.
\]
The link edge $u_{i+1}u_{i+2}$ is opposite the radial edge $vu_i$ in the
four-point configuration, so the three opposite pairs are equal.

It remains to exclude a flat configuration.  Any two consecutive vectors
$u_i-v$ and $u_{i+1}-v$ are linearly independent because they span a
nondegenerate triangular facet.  If the three vectors $u_i-v$ were linearly
dependent, all four points would lie in one plane $H$, and each of the three
distinct incident triangular facets would have affine hull $H$.  But a
facet plane of a convex three-polytope is supporting, and its intersection
with the polytope is one face.  The three facets would therefore coincide,
a contradiction.  Thus the four points form a nondegenerate tetrahedron.

Its three faces through $v$ are facets and hence copies of $T$.  Moreover,
the three radial labels in a closed degree-three link are precisely
$a,b,c$, and \eqref{eq:transition} permutes them to the three link labels.
Thus the fourth face $u_0u_1u_2$ also has side-length set $\{a,b,c\}$ and is
congruent to $T$.
\end{proof}

\begin{proof}[Proof of Theorem~\ref{thm:main}]
By Lemmas~\ref{lem:degree3} and \ref{lem:star}, the polytope contains four
vertices forming a nondegenerate tetrahedron with three equal pairs of
opposite edges and with every face congruent to $T$.

Call the vertices $A,B,C,D$ and translate their centroid to the origin.  Set
\[
 x=\frac{A+B-C-D}{4},\quad
 y=\frac{A-B+C-D}{4},\quad
 z=\frac{A-B-C+D}{4}.
\]
Then
\[
 A=x+y+z,\quad B=x-y-z,\quad
 C=-x+y-z,\quad D=-x-y+z.
\]
Equality of the opposite edges $AB$ and $CD$ gives
$|y+z|=|y-z|$, hence $y\cdot z=0$.  The other two opposite-pair equalities
similarly give $x\cdot z=x\cdot y=0$.  Nondegeneracy makes $x,y,z$ nonzero.

After permuting the coordinate names to match the three opposite-pair
lengths $a,b,c$, we have
\[
 a^2=4(|y|^2+|z|^2),\qquad
 b^2=4(|x|^2+|z|^2),\qquad
 c^2=4(|x|^2+|y|^2).
\]
Solving,
\begin{equation}\label{eq:disphenoid}
 |x|^2=\frac{b^2+c^2-a^2}{8},\quad
 |y|^2=\frac{c^2+a^2-b^2}{8},\quad
 |z|^2=\frac{a^2+b^2-c^2}{8}.
\end{equation}
All three quantities are strictly positive.  In particular,
$a^2+b^2-c^2>0$, so $T$ is acute.  This proves the first assertion and
contradicts $c^2\geq a^2+b^2$ in the stated right-or-obtuse range.
\end{proof}

\section{Verification and scope}

The proof above is deductive and does not depend on computation.  As a finite
cross-check, two short exact Python enumerators were used.  One independently
chooses any cyclic rotation on every incident face for each of the two
reference words $(a,b,c)$ and $(a,c,b)$.  For degrees $3,4,5,6,7,8$ it finds
respectively $3,0,0,3,0,0$ legal labeled stars in either handedness class,
and every degree-three output has the opposite-edge equalities used in
Lemma~\ref{lem:star}.  A second radial-word enumerator gives the same
divisibility pattern through degree $12$.  All comparisons are finite string
equalities; there is no floating-point tolerance or random search.

Exact finite enumeration with a separately replayable verifier is also used
in related discrete-geometric work of Zeng, Liu, Ratnavelu, and Huat
\cite{ZengEtAl2026Certificates}.  Here the enumerations are only checks of the
closed-form transition \eqref{eq:transition}, not theorem-bearing evidence.

Convexity is used in the supporting-plane argument of
Lemma~\ref{lem:star}; full dimensionality rules out the flat star;
triangularity supplies the cyclic vertex links and Euler count; scaleneness
makes the side transition deterministic; and the direct-copy hypothesis
selects one cyclic word rather than both.  The proof makes no assertion about
reflection-allowed realizations, acute constructions, isosceles triangles,
nonconvex surfaces, lower-dimensional configurations, polygonal facets, or
the classification of attainable face counts.

\paragraph{Computational and AI assistance.}
OpenAI Codex assisted with the literature search, proof organization,
construction and checking of the auxiliary exact enumerators, adversarial
review, and manuscript preparation.  A separate language-model review was
used as an additional error-finding pass, not as independent peer review.
Language-model output is not mathematical evidence: the complete argument is
given above, and the finite checks are reproducible exact computations.  No
proof assistant was used.

\bibliographystyle{plainurl}
\bibliography{references}

\end{document}
